\begin{thebibliography}{258}
\bibitem{1}\hypertarget{ref:1}{} M. Abouzaid, D. Auroux, A. I. Efimov, L. Katzarkov and D. Orlov, Homological mirror symmetry for punctured spheres. J. Amer. Math. Soc., 26 (2013), 1051–1083.
\bibitem{2}\hypertarget{ref:2}{} M. Abouzaid, D. Auroux and L. Katzarkov, Lagrangian fibrations on blowups of toric varieties and mirror symmetry for hypersurfaces. Publ. Math. Inst. Hautes Études Sci., 123 (2016), 199–282.
\bibitem{3}\hypertarget{ref:3}{} I. Assem, T. Brüstle, G. Charbonneau-Jodoin and P.-G. Plamondon, Gentle algebras arising from surface triangulations. Algebra Number Theory, 4 (2010), 201–229.
\bibitem{4}\hypertarget{ref:4}{} M. Abouzaid, On the Fukaya categories of higher genus surfaces. Adv. Math., 217 (2008), 1192–1235.
\bibitem{5}\hypertarget{ref:5}{} M. Abouzaid, Morse homology, tropical geometry, and homological mirror symmetry for toric varieties. Selecta Math. (N.S.), 15 (2009), 189–270.
\bibitem{6}\hypertarget{ref:6}{} M. Abouzaid, A geometric criterion for generating the Fukaya category. Publ. Math. Inst. Hautes Études Sci., 112 (2010), 191–240.
\bibitem{7}\hypertarget{ref:7}{} M. Abouzaid, A cotangent fibre generates the Fukaya category. Adv. Math., 228 (2011), 894–939.
\bibitem{8}\hypertarget{ref:8}{} L. Abrams, Two-dimensional topological quantum field theories and Frobenius algebras. J. Knot Theory Ramifications, 5 (1996), 569–587.
\bibitem{9}\hypertarget{ref:9}{} F. W. Anderson and K. R. Fuller, Rings and Categories of Modules, Graduate Texts in Mathematics, vol. 13. Springer, New York-Heidelberg, 1974.
\bibitem{10}\hypertarget{ref:10}{} L. Angeleri Hügel, D. Happel and H. Krause, eds, Handbook of Tilting Theory, London Mathematical Society Lecture Note Series, vol. 332. Cambridge University Press, Cambridge, 2007.
\bibitem{11}\hypertarget{ref:11}{} D. Auroux, L. Katzarkov and D. Orlov, Mirror symmetry for del Pezzo surfaces: vanishing cycles and coherent sheaves. Invent. Math., 166 (2006), 537–582.
\bibitem{12}\hypertarget{ref:12}{} V. I. Arnol’d, On a characteristic class entering into conditions of quantization. Funkcional. Anal. i Priložen., 1 (1967), 1–14.
\bibitem{13}\hypertarget{ref:13}{} M. Auslander, I. Reiten and S. O. Smalø, Representation Theory of Artin Algebras, Cambridge Studies in Advanced Mathematics, vol. 36. Cambridge University Press, Cambridge, 1997. Corrected reprint of the 1995 original.
\bibitem{14}\hypertarget{ref:14}{} I. Assem and A. Skowroński, Iterated tilted algebras of type $\widetilde{A}_n$. Math. Z., 195 (1987), 269–290.
\bibitem{15}\hypertarget{ref:15}{} M. Abouzaid and P. Seidel, An open string analogue of Viterbo functoriality. Geom. Topol., 14 (2010), 627–718.
\bibitem{16}\hypertarget{ref:16}{} P. S. Aspinwall, D-branes on Calabi-Yau manifolds. In Progress in String Theory. World Sci. Publ., Hackensack, NJ, 2005, pp. 1–152.
\bibitem{17}\hypertarget{ref:17}{} M. Artin, J. Tate and M. Van den Bergh, Some algebras associated to automorphisms of elliptic curves. In The Grothendieck Festschrift, Vol. I, Progress in Mathematics, vol. 86. Birkhäuser Boston, Boston, MA, 1990, pp. 33–85.
\bibitem{18}\hypertarget{ref:18}{} M. Atiyah, Topological quantum field theories. Publ. Math., Inst. Hautes Étud. Sci., 68 (1988), 175–186.
\bibitem{19}\hypertarget{ref:19}{} D. Auroux, Mirror symmetry and T -duality in the complement of an anticanonical divisor. J. Gökova Geom. Topol. GGT, 1 (2007), 51–91.
\bibitem{20}\hypertarget{ref:20}{} D. Auroux, A beginner’s introduction to Fukaya categories. In Contact and Symplectic Topology, Bolyai Society Mathematical Studies, vol. 26. János Bolyai Math. Soc., Budapest, 2014, pp. 85–136.
\bibitem{21}\hypertarget{ref:21}{} M. J. Bardzell, The alternating syzygy behavior of monomial algebras. J. Algebra, 188 (1997), 69–89.
\bibitem{22}\hypertarget{ref:22}{} V. V. Batyrev, Dual polyhedra and mirror symmetry for Calabi-Yau hypersurfaces in toric varieties. J. Algebraic Geom., 3 (1994), 493–535.
\bibitem{23}\hypertarget{ref:23}{} V. V. Batyrev and L. A. Borisov, Mirror duality and string-theoretic Hodge numbers. Invent. Math., 126 (1996), 183–203.
\bibitem{24}\hypertarget{ref:24}{} V. V. Batyrev and L. A. Borisov, On Calabi-Yau complete intersections in toric varieties. In Higher-Dimensional Complex Varieties (Trento, 1994). De Gruyter, Berlin, 1996, pp. 39–65.
\bibitem{25}\hypertarget{ref:25}{} V. V. Batyrev, I. Ciocan-Fontanine, B. Kim and D. van Straten, Conifold transitions and mirror symmetry for Calabi-Yau complete intersections in Grassmannians. Nuclear Phys. B, 514 (1998), 640–666.
\bibitem{26}\hypertarget{ref:26}{} A. A. Beilinson, The derived category of coherent sheaves on Pn . Selecta Math. Soviet., 3 (1983/84), 233–237. Selected translations.
\bibitem{27}\hypertarget{ref:27}{} A. Beilinson, V. Ginzburg and W. Soergel, Koszul duality patterns in representation theory. J. Amer. Math. Soc., 9 (1996), 473–527.
\bibitem{28}\hypertarget{ref:28}{} A. Banyaga and D. Hurtubise, Lectures on Morse Homology, Kluwer Texts in the Mathematical Sciences, vol. 29. Kluwer Academic Publishers Group, Dordrecht, 2004.
\bibitem{29}\hypertarget{ref:29}{} M. Bender and S. Mozgovoy, Crepant resolutions and brane tilings II: Tilting bundles. arXiv:0909.2013, 2009.
\bibitem{30}\hypertarget{ref:30}{} A. Bayer, E. Macrı̀and P. Stellari, The space of stability conditions on abelian threefolds, and on some Calabi-Yau threefolds. Invent. Math., 206 (2016), 869–933.
\bibitem{31}\hypertarget{ref:31}{} M. Bökstedt and A. Neeman, Homotopy limits in triangulated categories. Comp. Math., 86 (1993), 209–234.
\bibitem{32}\hypertarget{ref:32}{} A. Bondal and D. Orlov, Reconstruction of a variety from the derived category and groups of autoequivalences. Comp. Math., 125 (2001), 327–344.
\bibitem{33}\hypertarget{ref:33}{} R. Bocklandt, Consistency conditions for dimer models. Glasg. Math. J., 54 (2012), 429–447.
\bibitem{34}\hypertarget{ref:34}{} R. Bocklandt, Generating toric noncommutative crepant resolutions. J. Algebra, 364 (2012), 119–147.
\bibitem{35}\hypertarget{ref:35}{} R. Bocklandt, A dimer ABC. Bull. Lond. Math. Soc., 48 (2016), 387–451.
\bibitem{36}\hypertarget{ref:36}{} R. Bocklandt, Noncommutative mirror symmetry for punctured surfaces. Trans. Amer. Math. Soc., 368 (2016), 429–469. With an appendix by Mohammed Abouzaid.
\bibitem{37}\hypertarget{ref:37}{} R. Bocklandt, Strebel differentials and stable matrix factorizations. arXiv:1612.06800, 2016.
\bibitem{38}\hypertarget{ref:38}{} A. Bondal, Derived categories of toric varieties. In Convex and Algebraic Geometry, Oberwolfach Conference Reports, vol. 3. EMS Publishing House, Zürich, 2006, pp. 284–286.
\bibitem{39}\hypertarget{ref:39}{} T. Bridgeland, Stability conditions on triangulated categories. Ann. of Math. (2), 166 (2007), 317–345.
\bibitem{40}\hypertarget{ref:40}{} T. Bridgeland, Stability conditions on K3 surfaces. Duke Math. J., 141 (2008), 241–291.
\bibitem{41}\hypertarget{ref:41}{} M. Brion, Representations of quivers. In Geometric Methods in Representation Theory. I, Séminaires et Congrès, vol. 24. Soc. Math. France, Paris, 2012, pp. 103–144.
\bibitem{42}\hypertarget{ref:42}{} N. Broomhead, Dimer models and Calabi-Yau algebras. Mem. Amer. Math. Soc., 211 (2012).
\bibitem{43}\hypertarget{ref:43}{} P. Balmer and M. Schlichting, Idempotent completion of triangulated categories. J. Algebra, 236 (2001), 819–834.
\bibitem{44}\hypertarget{ref:44}{} T. Bridgeland and I. Smith, Quadratic differentials as stability conditions. Publ. Math. Inst. Hautes Études Sci., 121 (2015), 155–278.
\bibitem{45}\hypertarget{ref:45}{} R. Bott and L. W. Tu, Differential Forms in Algebraic Topology, Graduate Texts in Mathematics, vol. 82. Springer, New York-Berlin, 1982.
\bibitem{46}\hypertarget{ref:46}{} R.-O. Buchweitz, Maximal Cohen-Macaulay modules and Tate-cohomology over Gorenstein rings. University of Hannover, 1987. https://tspace .library.utoronto.ca/bitstream/1807/16682/1/maximal cohenmacaulay modules 1986.pdf
\bibitem{47}\hypertarget{ref:47}{} W. Crawley-Boevey, More lectures on representations of quivers. Preprint, University of Leeds, 1992. http://ambio1.leeds.ac.uk/pure/staff/ crawley b/morequivlecs.pdf
\bibitem{48}\hypertarget{ref:48}{} W. Crawley-Boevey, DMV lectures on representations of quivers, preprojective algebras and deformations of quotient singularities. Preprint, University of Leeds, 1999. www.math.uni-bielefeld.de/˜wcrawley/dmvlecs.pdf
\bibitem{49}\hypertarget{ref:49}{} T. Coates, A. Corti, S. Galkin, V. Golyshev and A. Kasprzyk, Mirror symmetry and Fano manifolds. In European Congress of Mathematics. Eur. Math. Soc., Zürich, 2013, pp. 285–300.
\bibitem{50}\hypertarget{ref:50}{} A. Caldararu, K. Costello and J. Tu, Categorical enumerative invariants, I: String vertices. arXiv:2009.06673, 2020.
\bibitem{51}\hypertarget{ref:51}{} P. Candelas, X. C. de la Ossa, P. S. Green and L. Parkes, A pair of Calabi-Yau manifolds as an exactly soluble superconformal theory. Nuclear Phys. B, 359 (1991), 21–74.
\bibitem{52}\hypertarget{ref:52}{} A. Cannas da Silva, Lectures on Symplectic Geometry, Lecture Notes in Mathematics, vol. 1764. Springer, Berlin, 2001.
\bibitem{53}\hypertarget{ref:53}{} K. Cieliebak, A. Floer and H. Hofer, Symplectic homology. II. A general construction. Math. Z., 218 (1995), 103–122.
\bibitem{54}\hypertarget{ref:54}{} B. Chantraine, An introduction to Fukaya categories. Preprint, University of Nantes, 2014. \url{www.math.sciences.univ-nantes.fr/~chantraine-b/en/FukayaCategory.pdf}
\bibitem{55}\hypertarget{ref:55}{} C.-H. Cho, H. Hong and S.-C. Lau, Localized mirror functor for Lagrangian immersions, and homological mirror symmetry for P1a,b,c . J. Differential Geom., 106 (2017), 45–126.
\bibitem{56}\hypertarget{ref:56}{} D. A. Cox and S. Katz, Mirror Symmetry and Algebraic Geometry, Mathematical Surveys and Monographs, vol. 68. American Mathematical Society, Providence, RI, 1999.
\bibitem{57}\hypertarget{ref:57}{} D. A. Cox, J. B. Little and H. K. Schenck, Toric Varieties, Graduate Studies in Mathematics, vol. 124. American Mathematical Society, Providence, RI, 2011.
\bibitem{58}\hypertarget{ref:58}{} K. Costello, Topological conformal field theories and Calabi-Yau categories. Adv. Math., 210 (2007), 165–214.
\bibitem{59}\hypertarget{ref:59}{} D. A. Cox, Mirror symmetry and polar duality of polytopes. Symmetry, 7 (2015), 1633–1645.
\bibitem{60}\hypertarget{ref:60}{} A. Căldăraru and J. Tu, Curved A∞ algebras and Landau-Ginzburg models. New York J. Math., 19 (2013), 305–342.
\bibitem{61}\hypertarget{ref:61}{} A. Caldararu and J. Tu, Computing a categorical Gromov-Witten invariant. arXiv:1706.09912, 2019.
\bibitem{62}\hypertarget{ref:62}{} A. Collinucci and T. Wyder, Introduction to topological string theory. Report No. KUL-TF-07/24, University of Leuven, 2007.
\bibitem{63}\hypertarget{ref:63}{} B. Davison, Consistency conditions for brane tilings. J. Algebra, 338 (2011), 1–23.
\bibitem{64}\hypertarget{ref:64}{} P. Deligne, P. Etingof, D. S. Freed, L. C. Jeffrey, D. Kazhdan, J. W. Morgan, D. R. Morrison and E. Witten, eds, Quantum Fields and Strings: A Course for Mathematicians. Vols. 1, 2. American Mathematical Society, Providence, RI; Institute for Advanced Study (IAS), Princeton, NJ, 1999. Material from the Special Year on Quantum Field Theory held at the Institute for Advanced Study, Princeton, NJ, 1996–1997.
\bibitem{65}\hypertarget{ref:65}{} T. Dyckerhoff and M. Kapranov, Triangulated surfaces in triangulated categories. J. Eur. Math. Soc. (JEMS), 20 (2018), 1473–1524.
\bibitem{66}\hypertarget{ref:66}{} P. Dolbeault, Sur la cohomologie des variétés analytiques complexes. C. R. Acad. Sci. Paris, 236 (1953), 175–177.
\bibitem{67}\hypertarget{ref:67}{} V. Drinfeld, DG quotients of DG categories. J. Algebra, 272 (2004), 643–91.
\bibitem{68}\hypertarget{ref:68}{} D. Dugger, A primer on homotopy colimits. Preprint, University of Oregon, 2008. https://pages.uoregon.edu/ddugger/hocolim.pdf
\bibitem{69}\hypertarget{ref:69}{} T. Dyckerhoff, Compact generators in categories of matrix factorizations. Duke Math. J., 159 (2011), 223–274.
\bibitem{70}\hypertarget{ref:70}{} A. I. Efimov, A proof of the Kontsevich-Soibelman conjecture. Mat. Sb., 202 (2011), 65–84.
\bibitem{71}\hypertarget{ref:71}{} A. I. Efimov, Homological mirror symmetry for curves of higher genus. Adv. Math., 230 (2012), 493–530.
\bibitem{72}\hypertarget{ref:72}{} D. Eisenbud and J. Harris, The Geometry of Schemes, Graduate Texts in Mathematics, vol. 197. Springer, New York, 2000.
\bibitem{73}\hypertarget{ref:73}{} D. Eisenbud, Homological algebra on a complete intersection, with an application to group representations. Trans. Amer. Math. Soc., 260 (1980), 35–64.
\bibitem{74}\hypertarget{ref:74}{} D. Eisenbud, Commutative Algebra, Graduate Texts in Mathematics, vol. 150. Springer, New York, 1995. With a view toward algebraic geometry.
\bibitem{75}\hypertarget{ref:75}{} G. Ellingsrud and S. A. Strømme, Bott’s formula and enumerative geometry. J. Amer. Math. Soc., 9 (1996), 175–193.
\bibitem{76}\hypertarget{ref:76}{} B. Fantechi, Stacks for everybody. In European Congress of Mathematics, Vol. I (Barcelona, 2000), Progress in Mathematics, vol. 201. Birkhäuser, Basel, 2001, pp. 349–359.
\bibitem{77}\hypertarget{ref:77}{} B. Fang, Homological mirror symmetry is T -duality for Pn . Commun. Number Theory Phys., 2 (2008), 719–742.
\bibitem{78}\hypertarget{ref:78}{} A. Floer, H. Hofer and K. Wysocki, Applications of symplectic homology. I. Math. Z., 217 (1994), 577–606.
\bibitem{79}\hypertarget{ref:79}{} A. Floer, Morse theory for Lagrangian intersections. J. Differential Geom., 28 (1988), 513–547.
\bibitem{80}\hypertarget{ref:80}{} B. Fang, C.-C. M. Liu, D. Treumann and E. Zaslow, T-duality and equivariant homological mirror symmetry for toric varieties. arXiv:0811.1228, 2008.
\bibitem{81}\hypertarget{ref:81}{} B. Fang, C.-C. M. Liu, D. Treumann and E. Zaslow, The coherent-constructible correspondence and homological mirror symmetry for toric varieties. In Geometry and Analysis. No. 2, Advanced Lectures in Mathematics, vol. 18. Int. Press, Somerville, MA, 2011, pp. 3–37.
\bibitem{82}\hypertarget{ref:82}{} K. Fukaya, Y.-G. Oh, H. Ohta and K. Ono, Canonical models of filtered A∞ -algebras and Morse complexes. In New Perspectives and Challenges in Symplectic Field Theory, CRM Proceedings and Lecture Notes, vol. 49. American Mathematical Society, Providence, RI, 2009, pp. 201–227.
\bibitem{83}\hypertarget{ref:83}{} K. Fukaya, Y.-G. Oh, H. Ohta and K. Ono, Lagrangian Intersection Floer Theory: Anomaly and Obstruction. Part I, AMS/IP Studies in Advanced Mathematics, vol. 46. American Mathematical Society, Providence, RI; International Press, Somerville, MA, 2009.
\bibitem{84}\hypertarget{ref:84}{} K. Fukaya, Y.-G. Oh, H. Ohta and K. Ono, Lagrangian Intersection Floer Theory: Anomaly and Obstruction. Part II, AMS/IP Studies in Advanced Mathematics, vol. 46. American Mathematical Society, Providence, RI; International Press, Somerville, MA, 2009.
\bibitem{85}\hypertarget{ref:85}{} W. Fulton and R. Pandharipande, Notes on stable maps and quantum cohomology. In Algebraic Geometry—Santa Cruz 1995, Proceedings of Symposia in Pure Mathematics, vol. 62. American Mathematical Society, Providence, RI, 1997, pp. 45–96.
\bibitem{86}\hypertarget{ref:86}{} M. Forsberg, M. Passare and A. Tsikh, Laurent determinants and arrangements of hyperplane amoebas. Adv. Math., 151 (2000), 45–70.
\bibitem{87}\hypertarget{ref:87}{} G. Friedman, Survey article: An elementary illustrated introduction to simplicial sets. Rocky Mountain J. Math., 42 (2012), 353–423.
\bibitem{88}\hypertarget{ref:88}{} K. R. Fuller, Biserial rings. In Ring Theory (Proc. Conf., Univ. Waterloo, Waterloo, 1978), Lecture Notes in Mathematics, vol. 734. Springer, Berlin, 1979, pp. 64–90.
\bibitem{89}\hypertarget{ref:89}{} W. Fulton, Introduction to Toric Varieties, Annals of Mathematics Studies, vol. 131. Princeton University Press, Princeton, NJ, 1993. The William H. Roever Lectures in Geometry.
\bibitem{90}\hypertarget{ref:90}{} E. Getzler, Batalin-Vilkovisky algebras and two-dimensional topological field theories. Comm. Math. Phys., 159 (1994), 265–285.
\bibitem{91}\hypertarget{ref:91}{} P. Griffiths and J. Harris, Principles of Algebraic Geometry. Wiley Classics Library. John Wiley \& Sons, New York, 1994. Reprint of the 1978 original.
\bibitem{92}\hypertarget{ref:92}{} A. B. Givental, Equivariant Gromov-Witten invariants. Int. Math. Res. Not. IMRN, 2015 (1996), 613–663.
\bibitem{93}\hypertarget{ref:93}{} E. Getzler and J. D. S. Jones, A∞ -algebras and the cyclic bar complex. Illinois J. Math., 34 (1990), 256–283.
\bibitem{94}\hypertarget{ref:94}{} B. Gammage and D. Nadler, Mirror symmetry for honeycombs. Trans. Amer. Math. Soc., 373 (2020), 71–107.
\bibitem{95}\hypertarget{ref:95}{} S. Ganatra, J. Pardon and V. Shende, Microlocal Morse theory of wrapped Fukaya categories. arXiv:1809.08807, 2018.
\bibitem{96}\hypertarget{ref:96}{} S. Ganatra, J. Pardon and V. Shende, Structural results in wrapped Floer theory. arXiv:1809.03427, 2018.
\bibitem{97}\hypertarget{ref:97}{} P. Gabriel and A. V. Roı̆ter, Representations of Finite-Dimensional Algebras, Encyclopaedia of Mathematical Sciences, vol. 73. Springer, Berlin, 1992. With a chapter by B. Keller.
\bibitem{98}\hypertarget{ref:98}{} A. Grothendieck, Sur la classification des fibrés holomorphes sur la sphère de Riemann. Amer. J. Math., 79 (1957), 121–138.
\bibitem{99}\hypertarget{ref:99}{} A. Grothendieck, Groupes de classes des catégories abéliennes et triangulées. Complexes parfaits. In Séminaire de Géométrie Algébrique du Bois-Marie 1965–66 SGA 5. Springer, 1977, pp. 351–371.
\bibitem{100}\hypertarget{ref:100}{} M. Gromov, Pseudo holomorphic curves in symplectic manifolds. Invent. Math., 82 (1985), 307–347.
\bibitem{101}\hypertarget{ref:101}{} M. Gross, Special Lagrangian fibrations. I. Topology. In Integrable Systems and Algebraic Geometry (Kobe/Kyoto, 1997). World Sci. Publ., River Edge, NJ, 1998, pp. 156–193.
\bibitem{102}\hypertarget{ref:102}{} M. Gross, Examples of special Lagrangian fibrations. In Symplectic Geometry and Mirror Symmetry (Seoul, 2000). World Sci. Publ., River Edge, NJ, 2001, pp. 81–109.
\bibitem{103}\hypertarget{ref:103}{} M. Gross, Special Lagrangian fibrations. II. Geometry. A survey of techniques in the study of special Lagrangian fibrations. In Winter School on Mirror Symmetry, Vector Bundles and Lagrangian Submanifolds (Cambridge, MA, 1999), AMS/IP Studies in Advanced Mathematics, vol. 23. American Mathematical Society, Providence, RI, 2001, pp. 95–150.
\bibitem{104}\hypertarget{ref:104}{} M. Gross and B. Siebert, Mirror symmetry via logarithmic degeneration data. I. J. Differential Geom., 72 (2006), 169–338.
\bibitem{105}\hypertarget{ref:105}{} M. Gross and B. Siebert, Mirror symmetry via logarithmic degeneration data, II. J. Algebraic Geom., 19 (2010), 679–780.
\bibitem{106}\hypertarget{ref:106}{} I. M. Gelfand, A. V. Zelevinskiı̆ and M. M. Kapranov, A-discriminants and Cayley-Koszul complexes. Dokl. Akad. Nauk SSSR, 307 (1989), 1307–1311.
\bibitem{107}\hypertarget{ref:107}{} M. Habermann, Homological mirror symmetry for invertible polynomials in two variables. arXiv:2003.01106, 2020.
\bibitem{108}\hypertarget{ref:108}{} D. Happel, Triangulated Categories in the Representation Theory of Finite-Dimensional Algebras, London Mathematical Society Lecture Note Series, vol. 119. Cambridge University Press, Cambridge, 1988.
\bibitem{109}\hypertarget{ref:109}{} R. Hartshorne, Algebraic Geometry, Graduate Texts in Mathematics, vol. 52. Springer, New York-Heidelberg, 1977.
\bibitem{110}\hypertarget{ref:110}{} A. Hatcher, On triangulations of surfaces. Topology Appl., 40 (1991), 189–194.
\bibitem{111}\hypertarget{ref:111}{} A. Hatcher, Algebraic Topology. Cambridge University Press, Cambridge, 2002.
\bibitem{112}\hypertarget{ref:112}{} L. Hille and J. A. de la Peña, Stable representations of quivers. J. Pure Appl. Algebra, 172 (2002), 205–224.
\bibitem{113}\hypertarget{ref:113}{} N. Hitchin, Lectures on special Lagrangian submanifolds. In Winter School on Mirror Symmetry, Vector Bundles and Lagrangian Submanifolds (Cambridge, MA, 1999), AMS/IP Studies in Advanced Mathematics, vol. 23. American Mathematical Society, Providence, RI, 2001, pp. 151–182.
\bibitem{114}\hypertarget{ref:114}{} K. Hori, S. Katz, A. Klemm, R. Pandharipande, R. Thomas, C. Vafa, R. Vakil and E. Zaslow, Mirror Symmetry, Clay Mathematics Monographs, vol. 1. American Mathematical Society, Providence, RI; Clay Mathematics Institute, Cambridge, MA, 2003. With a preface by Vafa.
\bibitem{115}\hypertarget{ref:115}{} F. Haiden, L. Katzarkov and M. Kontsevich, Flat surfaces and stability structures. Publ. Math. Inst. Hautes Études Sci., 126 (2017), 247–318.
\bibitem{116}\hypertarget{ref:116}{} G. Hochschild, B. Kostant and A. Rosenberg, Differential forms on regular affine algebras. Trans. Amer. Math. Soc., 102 (1962), 383–408.
\bibitem{117}\hypertarget{ref:117}{} D. Huybrechts and M. Lehn, The Geometry of Moduli Spaces of Sheaves, 2nd ed., Cambridge Mathematical Library. Cambridge University Press, Cambridge, 2010.
\bibitem{118}\hypertarget{ref:118}{} M. Hazewinkel and C. F Martin, A short elementary proof of Grothendieck’s theorem on algebraic vectorbundles over the projective line. J. Pure Appl. Algebra, 25 (1982), 207–211.
\bibitem{119}\hypertarget{ref:119}{} G. Hochschild, On the cohomology groups of an associative algebra. Ann. of Math. (2), 46 (1945), 58–67.
\bibitem{120}\hypertarget{ref:120}{} L. Hille and M. Perling, Tilting bundles on rational surfaces and quasi-hereditary algebras. Ann. Inst. Fourier (Grenoble), 64 (2014), 625–644.
\bibitem{121}\hypertarget{ref:121}{} A. Hanany and R.-K. Seong, Brane tilings and specular duality. J. High Energy Phys., 2012 (2012), 107.
\bibitem{122}\hypertarget{ref:122}{} M. Henneaux and C. Teitelboim, BRST cohomology in classical mechanics. Comm. Math. Phys., 115 (1988), 213–230.
\bibitem{123}\hypertarget{ref:123}{} M. Hutchings, Lecture notes on Morse homology (with an eye towards Floer theory and pseudoholomorphic curves). Preprint, Berkeley, 2002. https:// math.berkeley.edu/˜hutching/teach/276-2010/mfp.ps
\bibitem{124}\hypertarget{ref:124}{} D. Huybrechts, Fourier-Mukai Transforms in Algebraic Geometry. Oxford Mathematical Monographs. Clarendon Press, Oxford, 2006.
\bibitem{125}\hypertarget{ref:125}{} A. Ishii and K. Ueda, On moduli spaces of quiver representations associated with dimer models. In Higher Dimensional Algebraic Varieties and Vector Bundles, RIMS Kôkyûroku Bessatsu, B9. Res. Inst. Math. Sci. (RIMS), Kyoto, 2008, pp. 127–141.
\bibitem{126}\hypertarget{ref:126}{} A. Ishii and K. Ueda, A note on consistency conditions on dimer models. In Higher Dimensional Algebraic Geometry, RIMS Kôkyûroku Bessatsu, B24. Res. Inst. Math. Sci. (RIMS), Kyoto, 2011, pp. 143–164.
\bibitem{127}\hypertarget{ref:127}{} A. Ishii and K. Ueda, Dimer models and the special McKay correspondence. Geom. Topol., 19 (2015), 3405–3466.
\bibitem{128}\hypertarget{ref:128}{} D. Joyce, Lectures on Calabi-Yau and special Lagrangian geometry. arXiv:math/0108088, 2001.
\bibitem{129}\hypertarget{ref:129}{} D. Joyce, Singularities of special Lagrangian fibrations and the SYZ conjecture. Comm. Anal. Geom., 11 (2003), 859–907.
\bibitem{130}\hypertarget{ref:130}{} D. D. Joyce, Riemannian Holonomy Groups and Calibrated Geometry, Oxford Graduate Texts in Mathematics, vol. 12. Oxford University Press, Oxford, 2007.
\bibitem{131}\hypertarget{ref:131}{} T. V. Kadeishvili, The algebraic structure in the homology of an A(∞)-algebra. Soobshch. Akad. Nauk Gruzin. SSR, 108 (1982, 1983), 249–252.
\bibitem{132}\hypertarget{ref:132}{} P. W. Kasteleyn, Dimer statistics and phase transitions. J. Mathematical Phys., 4 (1963), 287–293.
\bibitem{133}\hypertarget{ref:133}{} B. Keller, Introduction to A-infinity algebras and modules. Homology Homotopy Appl., 3 (2001), 1–35.
\bibitem{134}\hypertarget{ref:134}{} B. Keller, Derived invariance of higher structures on the Hochschild complex. Preprint, 2003. \url{https://webusers.imj-prg.fr/~bernhard.keller/publ/dih.pdf}
\bibitem{135}\hypertarget{ref:135}{} B. Keller, A-infinity algebras, modules and functor categories. In Trends in Representation Theory of Algebras and Related Topics, Contemporary Mathematics, vol. 406. American Mathematical Society, Providence, RI, 2006, pp. 67–93.
\bibitem{136}\hypertarget{ref:136}{} B. Keller, Derived categories and tilting. In Handbook of Tilting Theory, London Mathematical Society Lecture Note Series, vol. 332. Cambridge University Press, Cambridge, 2007, pp. 49–104.
\bibitem{137}\hypertarget{ref:137}{} R. Kenyon, An introduction to the dimer model. In School and Conference on Probability Theory, ICTP Lecture Notes, XVII. Abdus Salam Int. Cent. Theoret. Phys., Trieste, 2004, pp. 267–304.
\bibitem{138}\hypertarget{ref:138}{} A. D. King, Moduli of representations of finite-dimensional algebras. Quart. J. Math. Oxford Ser. (2), 45 (1994), 515–530.
\bibitem{139}\hypertarget{ref:139}{} A. Kapustin, L. Katzarkov, D. Orlov and M. Yotov, Homological mirror symmetry for manifolds of general type. Cent. Eur. J. Math., 7 (2009), 571–605.
\bibitem{140}\hypertarget{ref:140}{} H. Knörrer, Cohen-Macaulay modules on hypersurface singularities I. Invent. Math., 88 (1987), 153–164.
\bibitem{141}\hypertarget{ref:141}{} M. Kontsevich, Homological algebra of mirror symmetry. In Proceedings of the International Congress of Mathematicians, Vol. 1, 2 (Zürich, 1994). Birkhäuser, Basel, 1995, pp. 120–139.
\bibitem{142}\hypertarget{ref:142}{} M. Kontsevich, Triangulated categories and geometry. Course at the École Normale Supérieure, Paris, notes taken by J. Bellaıche, 1998.
\bibitem{143}\hypertarget{ref:143}{} M. Kontsevich, Deformation quantization of Poisson manifolds. Lett. Math. Phys., 66 (2003), 157–216.
\bibitem{144}\hypertarget{ref:144}{} M. Kontsevich, Symplectic geometry of homological algebra. Preprint, IHES, 2009. www.ihes.fr/˜/maxim/TEXTS/Symplectic\_AT2009.pdf
\bibitem{145}\hypertarget{ref:145}{} M. Kontsevich and A. L. Rosenberg, Noncommutative smooth spaces. In The Gelfand Mathematical Seminars, 1996–1999, Gelfand Math. Sem. Birkhäuser Boston, Boston, MA, 2000, pp. 85–108.
\bibitem{146}\hypertarget{ref:146}{} M. Kontsevich and Y. Soibelman, Homological mirror symmetry and torus fibrations. In Symplectic Geometry and Mirror Symmetry (Seoul, 2000). World Sci. Publ., River Edge, NJ, 2001, pp. 203–263.
\bibitem{147}\hypertarget{ref:147}{} M. Kontsevich and Y. Soibelman, Deformation theory. Livre en préparation, 2002. https://people.maths.ox.ac.uk/beem/papers/kontsevich soibelman deformation theory 1.pdf
\bibitem{148}\hypertarget{ref:148}{} M. Kashiwara and P. Schapira, Categories and sheaves. Grundlehren der Mathematischen Wissenschaften [Fundamental Principles of Mathematical Sciences], vol. 332. Springer, Berlin, 2006.
\bibitem{149}\hypertarget{ref:149}{} M. Kontsevich and Y. Soibelman, Stability structures, motivic DonaldsonThomas invariants and cluster transformations. arXiv:0811.2435, 2008.
\bibitem{150}\hypertarget{ref:150}{} M. Kontsevich and Y. Soibelman, Notes on A∞ -algebras, A∞ -categories and non-commutative geometry. In Homological Mirror Symmetry, Lecture Notes in Physics, vol. 757. Springer, Berlin, 2009, pp. 153–219.
\bibitem{151}\hypertarget{ref:151}{} T. Y. Lam, Lectures on Modules and Rings, Graduate Texts in Mathematics, vol. 189. Springer, New York, 1999.
\bibitem{152}\hypertarget{ref:152}{} L. Le Bruyn, Noncommutative Geometry and Cayley-Smooth Orders, Pure and Applied Mathematics (Boca Raton), vol. 290. Chapman \& Hall/CRC, Boca Raton, FL, 2008.
\bibitem{153}\hypertarget{ref:153}{} W. Lowen and M. Van den Bergh, The curvature problem for formal and infinitesimal deformations. arXiv:1505.03698, 2015.
\bibitem{154}\hypertarget{ref:154}{} S. Lau, C.-H. Cho and H. Hong, Noncommutative homological mirror functor. Mem. Amer. Math. Soc., 2019.
\bibitem{155}\hypertarget{ref:155}{} H. M. Lee, Homological Mirror Symmetry for Open Riemann Surfaces from Pair-of-Pants Decompositions. ProQuest LLC, Ann Arbor, MI, 2015. Thesis (Ph.D.)–University of California, Berkeley.
\bibitem{156}\hypertarget{ref:156}{} D. Labardini-Fragoso, Quivers with potentials associated to triangulated surfaces. Proc. Lond. Math. Soc. (3), 98 (2009), 797–839.
\bibitem{157}\hypertarget{ref:157}{} J. Lipman, Notes on derived functors and Grothendieck duality. In Foundations of Grothendieck Duality for Diagrams of Schemes, Lecture Notes in Mathematics, vol. 1960. Springer, Berlin, 2009, pp. 1–259.
\bibitem{158}\hypertarget{ref:158}{} B. H. Lian, K. Liu and S.-T. Yau, Mirror principle. I. Asian J. Math., 1 (1997), 729–763.
\bibitem{159}\hypertarget{ref:159}{} V. A. Lunts and D. O. Orlov, Uniqueness of enhancement for triangulated categories. J. Amer. Math. Soc., 23 (2010), 853–908.
\bibitem{160}\hypertarget{ref:160}{} J.-L. Loday, Cyclic Homology, vol. 301. Springer, 2013.
\bibitem{161}\hypertarget{ref:161}{} Y. Lekili and T. Perutz, Fukaya categories of the torus and Dehn surgery. Proc. Natl. Acad. Sci. USA, 108 (2011), 8106–8113.
\bibitem{162}\hypertarget{ref:162}{} Y. Lekili and T. Perutz, Arithmetic mirror symmetry for the 2-torus. arXiv:1211.4632, 2012.
\bibitem{163}\hypertarget{ref:163}{} K. H. Lin and D. Pomerleano, Global matrix factorizations. Math. Res. Lett., 20 (2013), 91–106.
\bibitem{164}\hypertarget{ref:164}{} Y. Lekili and A. Polishchuk, Arithmetic mirror symmetry for genus 1 curves with n marked points. Selecta Math. (N.S.), 23 (2017), 1851–1907.
\bibitem{165}\hypertarget{ref:165}{} Y. Lekili and A. Polishchuk, Auslander orders over nodal stacky curves and partially wrapped Fukaya categories. J. Topol., 11 (2018), 615–644.
\bibitem{166}\hypertarget{ref:166}{} Y. Lekili and A. Polishchuk, Derived equivalences of gentle algebras via Fukaya categories. Math. Ann., 376 (2020), 187–225.
\bibitem{167}\hypertarget{ref:167}{} D.-M. Lu, J. H. Palmieri, Q.-S. Wu and J. J. Zhang, Koszul equivalences in A∞ -algebras. New York J. Math., 14 (2008), 325–378.
\bibitem{168}\hypertarget{ref:168}{} D.-M. Lu, J. H. Palmieri, Q.-S. Wu and J. J. Zhang, A-infinity structure on Extalgebras. J. Pure Appl. Algebra, 213 (2009), 2017–2037.
\bibitem{169}\hypertarget{ref:169}{} Y. Lekili and K. Ueda, Homological mirror symmetry for K3 surfaces via moduli of A∞ -structures. arXiv:1806.04345, 2018.
\bibitem{170}\hypertarget{ref:170}{} M. Markl, Transferring A∞ (strongly homotopy associative) structures. Rend. Circ. Mat. Palermo (2) Suppl., 79 (2006), 139–151.
\bibitem{171}\hypertarget{ref:171}{} W. S. Massey, Some higher order cohomology operations. In Symposium Internacional de Topologı́a Algebraica International Symposium on Algebraic Topology. Universidad Nacional Autónoma de México and UNESCO, Mexico City, 1958, pp. 145–154.
\bibitem{172}\hypertarget{ref:172}{} W. S. Massey, A Basic Course in Algebraic Topology, Graduate Texts in Mathematics, vol. 127. Springer, New York, 1991.
\bibitem{173}\hypertarget{ref:173}{} V. P Maslov, V. C. Bouslaev and V. I. Arnol’d, Théorie des perturbations et méthodes asymptotiques, Dunod, Paris, 1972.
\bibitem{174}\hypertarget{ref:174}{} S. A. Merkulov, Strong homotopy algebras of a Kähler manifold. Int. Math. Res. Not. IMRN, 1999 (1999), 153–164.
\bibitem{175}\hypertarget{ref:175}{} G. Mikhalkin, Amoebas of algebraic varieties and tropical geometry. In Different Faces of Geometry, International Mathematical Series (N. Y.), vol. 3. Kluwer/Plenum, New York, 2004, pp. 257–300.
\bibitem{176}\hypertarget{ref:176}{} S. Mac Lane, Categories for the Working Mathematician, 2nd ed., Graduate Texts in Mathematics, vol. 5. Springer, New York, 1998.
\bibitem{177}\hypertarget{ref:177}{} K. Morita, Duality for modules and its applications to the theory of rings with minimum condition. Sci. Rep. Tokyo Kyoiku Daigaku Sect. A, 6 (1958), 83–142.
\bibitem{178}\hypertarget{ref:178}{} I. Mori, S. Okawa and K. Ueda, Moduli of noncommutative Hirzebruch surfaces. arXiv:1903.06457, 2019.
\bibitem{179}\hypertarget{ref:179}{} S. Mozgovoy, Crepant resolutions and brane tilings I: Toric realization. arXiv:0908.3475, 2009.
\bibitem{180}\hypertarget{ref:180}{} S. Mozgovoy and M. Reineke, On the noncommutative Donaldson-Thomas invariants arising from brane tilings. Adv. Math., 223 (2010), 1521–1544.
\bibitem{181}\hypertarget{ref:181}{} D. Maclagan and B. Sturmfels, Introduction to Tropical Geometry, Graduate Studies in Mathematics, vol. 161. American Mathematical Society, Providence, RI, 2015.
\bibitem{182}\hypertarget{ref:182}{} D. McDuff and D. Salamon, Introduction to Symplectic Topology, 3rd ed., Oxford Graduate Texts in Mathematics. Oxford University Press, Oxford, 2017.
\bibitem{183}\hypertarget{ref:183}{} S. Mukai, Duality between D(X) and D(X̂) with its application to Picard sheaves. Nagoya Math. J., 81 (1981), 153–175.
\bibitem{184}\hypertarget{ref:184}{} D. Mumford, The Red Book of Varieties and Schemes, expanded ed., Lecture Notes in Mathematics, vol. 1358. Springer, Berlin, 1999. Includes the Michigan lectures (1974) on curves and their Jacobians. With contributions by Enrico Arbarello.
\bibitem{185}\hypertarget{ref:185}{} D. Murfet, Derived categories of quasi-coherent sheaves. Available at therisingsea.org, 2006.
\bibitem{186}\hypertarget{ref:186}{} D. Murfet, Derived categories of sheaves. Available at therisingsea.org, 2006.
\bibitem{187}\hypertarget{ref:187}{} D. Nadler, Microlocal branes are constructible sheaves. Selecta Math. (N.S.), 15 (2009), 563–619.
\bibitem{188}\hypertarget{ref:188}{} D. Nadler, Fukaya categories as categorical Morse homology. SIGMA Symmetry Integrability Geom. Methods Appl., 10 (2014) Paper 018, 47 pages.
\bibitem{189}\hypertarget{ref:189}{} D. Nadler, Cyclic symmetries of An -quiver representations. Adv. Math., 269 (2015), 346–363.
\bibitem{190}\hypertarget{ref:190}{} D. Nadler, Wrapped microlocal sheaves on pairs of pants. arXiv:1604.00114, 2016.
\bibitem{191}\hypertarget{ref:191}{} D. Nadler, Arboreal singularities. Geom. Topol., 21 (2017), 1231–1274.
\bibitem{192}\hypertarget{ref:192}{} L. A. Nazarova and A. V. Roiter, A problem of I. M. Gel’fand. Funct. Anal. Appl., 7 (1973), 299–311.
\bibitem{193}\hypertarget{ref:193}{} M. S. Narasimhan and C. S. Seshadri, Stable and unitary vector bundles on a compact Riemann surface. Ann. of Math. (2), 82 (1965), 540–567.
\bibitem{194}\hypertarget{ref:194}{} Y. Nohara and K. Ueda, Homological mirror symmetry for the quintic 3-fold. Geom. Topol., 16 (2012), 1967–2001.
\bibitem{195}\hypertarget{ref:195}{} D. Nadler and E. Zaslow, Constructible sheaves and the Fukaya category. J. Amer. Math. Soc., 22 (2009), 233–286.
\bibitem{196}\hypertarget{ref:196}{} M. Olsson, Algebraic Spaces and Stacks, American Mathematical Society Colloquium Publications, vol. 62. American Mathematical Society, Providence, RI, 2016.
\bibitem{197}\hypertarget{ref:197}{} S. Opper, P.-G. Plamondon and S. Schroll, A geometric model for the derived category of gentle algebras. arXiv:1801.09659, 2018.
\bibitem{198}\hypertarget{ref:198}{} D. O. Orlov, Derived categories of coherent sheaves and equivalences between them. Uspekhi Mat. Nauk, 58 (2003), 89–172.
\bibitem{199}\hypertarget{ref:199}{} D. O. Orlov, Triangulated categories of singularities, and equivalences between Landau-Ginzburg models. Mat. Sb., 197 (2006), 117–132.
\bibitem{200}\hypertarget{ref:200}{} D. Orlov, Derived categories of coherent sheaves and triangulated categories of singularities. In Algebra, Arithmetic, and Geometry: In Honor of Yu. I. Manin. Vol. II, Progress in Mathematics, vol. 270. Birkhäuser Boston, Boston, MA, 2009, pp. 503–531.
\bibitem{201}\hypertarget{ref:201}{} D. Orlov, Matrix factorizations for nonaffine LG-models. Math. Ann., 353 (2012), 95–108.
\bibitem{202}\hypertarget{ref:202}{} A. Polishchuk, A field guide to A-infinity sign conventions. Preprint. https:// pages.uoregon.edu/apolish/ainf-signs.pdf
\bibitem{203}\hypertarget{ref:203}{} J. P. Pridham, Unifying derived deformation theories. Adv. Math., 224 (2010), 772–826.
\bibitem{204}\hypertarget{ref:204}{} J. Pascaleff and N. Sibilla, Topological Fukaya category and mirror symmetry for punctured surfaces. Compos. Math., 155 (2019), 599–644.
\bibitem{205}\hypertarget{ref:205}{} A. Polishchuk and E. Zaslow, Categorical mirror symmetry: The elliptic curve. Adv. Theor. Math. Phys., 2 (1998), 443–470.
\bibitem{206}\hypertarget{ref:206}{} D. Quillen, Projective modules over polynomial rings. Invent. Math., 36 (1976), 167–171.
\bibitem{207}\hypertarget{ref:207}{} M. Reineke, Moduli of representations of quivers. In Trends in Representation Theory of Algebras and Related Topics, EMS Ser. Congr. Rep. European Mathematical Society, Zürich, 2008, pp. 589–637.
\bibitem{208}\hypertarget{ref:208}{} C. M. Ringel, The repetitive algebra of a gentle algebra, Bol. Soc. Mat. Mexicana, 3 (1997), 235–253.
\bibitem{209}\hypertarget{ref:209}{} C. P. Rourke and B. J. Sanderson, $\Delta$-sets. I. Homotopy theory. Quart. J. Math. Oxford Ser. (2), 22 (1971), 321–338.
\bibitem{210}\hypertarget{ref:210}{} A. Rizzardo and M. Van den Bergh, A note on non-unique enhancements. Proc. Amer. Math. Soc., 147 (2019), 451–453.
\bibitem{211}\hypertarget{ref:211}{} A. Rizzardo and M. Van den Bergh, A k-linear triangulated category without a model. Ann. of Math. (2), 191 (2020), 393–437.
\bibitem{212}\hypertarget{ref:212}{} D. Salamon, Lectures on Floer homology. In Symplectic Geometry and Topology, IAS/Park City Mathematics Series, vol. 7. American Mathematical Society, Providence, RI, 1999, pp. 143–229.
\bibitem{213}\hypertarget{ref:213}{} R. Schiffler, Quiver Representations. CMS Books in Mathematics/Ouvrages de Mathématiques de la SMC. Springer, Cham, 2014.
\bibitem{214}\hypertarget{ref:214}{} G. Segal, Lecture notes on topological field theory. Preprint, 1999. http:// web.math.ucsb.edu/˜drm/conferences/ITP99/segal/
\bibitem{215}\hypertarget{ref:215}{} G. Segal, The definition of conformal field theory. In Topology, Geometry and Quantum Field Theory, London Mathematical Society Lecture Note Series, vol. 308. Cambridge University Press, Cambridge, 2004, pp. 421–577.
\bibitem{216}\hypertarget{ref:216}{} E. Segal, Equivalence between GIT quotients of Landau-Ginzburg B-models. Comm. Math. Phys., 304 (2011), 411–432.
\bibitem{217}\hypertarget{ref:217}{} P. Seidel, Fukaya categories and deformations. In Proceedings of the International Congress of Mathematicians, Vol. II (Beijing, 2002). Higher Ed. Press, Beijing, 2002, pp. 351–360.
\bibitem{218}\hypertarget{ref:218}{} P. Seidel, Fukaya Categories and Picard-Lefschetz Theory. Zurich Lectures in Advanced Mathematics. European Mathematical Society (EMS), Zürich, 2008.
\bibitem{219}\hypertarget{ref:219}{} P. Seidel, Suspending Lefschetz fibrations, with an application to local mirror symmetry. Comm. Math. Phys., 297 (2010), 515–528.
\bibitem{220}\hypertarget{ref:220}{} P. Seidel, Homological mirror symmetry for the genus two curve. J. Algebraic Geom., 20 (2011), 727–769.
\bibitem{221}\hypertarget{ref:221}{} P. Seidel, Homological mirror symmetry for the quartic surface. Mem. Amer. Math. Soc., 236 (2015).
\bibitem{222}\hypertarget{ref:222}{} J.-P. Serre, Faisceaux algébriques cohérents. Ann. of Math. (2), 61 (1955), 197–278.
\bibitem{223}\hypertarget{ref:223}{} N. Sheridan, Homological mirror symmetry for Calabi-Yau hypersurfaces in projective space. Invent. Math., 199 (2015), 1–186.
\bibitem{224}\hypertarget{ref:224}{} N. Sheridan, Versality of the relative Fukaya category. Geom. Topol., 24 (2020), 747–884.
\bibitem{225}\hypertarget{ref:225}{} I. Shipman, A geometric approach to Orlov’s theorem. Compos. Math., 148 (2012), 1365–1389.
\bibitem{226}\hypertarget{ref:226}{} S. P. Smith, Non-commutative algebraic geometry. Preprint, Lecture Notes at University of Washington, 2000. \url{https://sites.math.washington.edu/~smith/Teaching/513nag/ch1.ps}
\bibitem{227}\hypertarget{ref:227}{} I. Smith, A symplectic prolegomenon. Bull. Amer. Math. Soc. (N.S.), 52 (2015), 415–464.
\bibitem{228}\hypertarget{ref:228}{} D. Sorensen, Classification of vector bundles over P1. Preprint. www.derekhsorensen.com/docs/sorensen-classification-vectorbundles.pdf
\bibitem{229}\hypertarget{ref:229}{} J. D. Stasheff, Homotopy associativity of H-spaces. I, II. Trans. Amer. Math. Soc. 108 (1963), 275–292, 293–312.
\bibitem{230}\hypertarget{ref:230}{} The Stacks Project Authors. Stacks Project. https://stacks.math .columbia.edu, 2018.
\bibitem{231}\hypertarget{ref:231}{} The structure of Coh $\mathbb P^1$. Preprint. \url{http://people.math.harvard.edu/~hirolee/pdfs/280x-13-14-dbcoh.pdf}
\bibitem{232}\hypertarget{ref:232}{} N. Sibilla, D. Treumann and E. Zaslow, Ribbon graphs and mirror symmetry. Selecta Math. (N.S.), 20 (2014), 979–1002.
\bibitem{233}\hypertarget{ref:233}{} T. Sutherland, The modular curve as the space of stability conditions of a CY3 algebra. arXiv:1111.4184, 2011.
\bibitem{234}\hypertarget{ref:234}{} A. Skowroński and J. Waschbüsch, Representation-finite biserial algebras. J. Reine Angew. Math., 345 (1983), 172–181.
\bibitem{235}\hypertarget{ref:235}{} R. G. Swan, Vector bundles and projective modules. Trans. Amer. Math. Soc., 105 (1962), 264–277.
\bibitem{236}\hypertarget{ref:236}{} Z. Sylvan, On partially wrapped Fukaya categories. J. Topol., 12 (2019), 372–441.
\bibitem{237}\hypertarget{ref:237}{} A. Strominger, S.-T. Yau and E. Zaslow, Mirror symmetry is T -duality. In Winter School on Mirror Symmetry, Vector Bundles and Lagrangian Submanifolds (Cambridge, MA, 1999), AMS/IP Studies in Advanced Mathematics, vol. 23. American Mathematical Society, Providence, RI, 2001, pp. 333–347.
\bibitem{238}\hypertarget{ref:238}{} D. Salamon and E. Zehnder, Morse theory for periodic solutions of Hamiltonian systems and the Maslov index. Comm. Pure Appl. Math., 45 (1992), 1303–1360.
\bibitem{239}\hypertarget{ref:239}{} G. Tabuada, Une structure de catégorie de modèles de Quillen sur la catégorie des dg-catégories. C. R. Math. Acad. Sci. Paris, 340 (2005), 15–19.
\bibitem{240}\hypertarget{ref:240}{} G. Tabuada, Finite generation of the numerical Grothendieck group. arXiv:1704.06252, 2017.
\bibitem{241}\hypertarget{ref:241}{} A. A. Takeda, Relative stability conditions on Fukaya categories of surfaces. arXiv:1811.10592, 2018.
\bibitem{242}\hypertarget{ref:242}{} D. E. Tamarkin, Operadic Proof of M. Kontsevich’s Formality Theorem. ProQuest LLC, Ann Arbor, MI, 1999. Thesis (Ph.D.)–The Pennsylvania State University.
\bibitem{243}\hypertarget{ref:243}{} B. Toen, The homotopy theory of dg-categories and derived Morita theory. Invent. Math., 167 (2007), 615–667.
\bibitem{244}\hypertarget{ref:244}{} B. Toen, Lectures on dg-categories. In Topics in Algebraic and Topological K-Theory, Lecture Notes in Mathematics, vol. 2008. Springer, Berlin, 2011, pp. 243–302.
\bibitem{245}\hypertarget{ref:245}{} J. Tu, Homological mirror symmetry and Fourier-Mukai transform. Int. Math. Res. Not. IMRN, 2015 (2015), 579–630.
\bibitem{246}\hypertarget{ref:246}{} V. Turaev and A. Virelizier, Monoidal Categories and Topological Field Theory, Progress in Mathematics, vol. 322. Birkhäuser/Springer, Cham, 2017.
\bibitem{247}\hypertarget{ref:247}{} K. Ueda, Homological mirror symmetry and simple elliptic singularities. arXiv math/0604361, 2006.
\bibitem{248}\hypertarget{ref:248}{} R. Vakil, The rising sea: Foundations of algebraic geometry. Preprint, 2017. \url{http://math.stanford.edu/~vakil/216blog/FOAGapr2915public.pdf}
\bibitem{249}\hypertarget{ref:249}{} M. Van den Bergh, Calabi-Yau algebras and superpotentials. Selecta Math. (N.S.), 21 (2015), 555–603.
\bibitem{250}\hypertarget{ref:250}{} M. van Garrel, D. P. Overholser and H. Ruddat, Enumerative aspects of the Gross-Siebert program. In Calabi-Yau Varieties: Arithmetic, Geometry and Physics. Fields Institute Monographs, vol. 34. Fields Inst. Res. Math. Sci., Toronto, ON, 2015, pp. 337–420.
\bibitem{251}\hypertarget{ref:251}{} M. Vonk, A mini-course on topological strings. arXiv:hep-th/0504147, 2005.
\bibitem{252}\hypertarget{ref:252}{} A. Weinstein, Symplectic manifolds and their Lagrangian submanifolds. Adv. Math., 6 (1971), 329–346.
\bibitem{253}\hypertarget{ref:253}{} C. A. Weibel, An Introduction to Homological Algebra, Cambridge Studies in Advanced Mathematics, vol. 38. Cambridge University Press, Cambridge, 1994.
\bibitem{254}\hypertarget{ref:254}{} C. A Weibel, The K-Book: An Introduction to Algebraic K-Theory, vol. 145. American Mathematical Society, Providence, RI, 2013.
\bibitem{255}\hypertarget{ref:255}{} C. Wendl, A beginner’s overview of symplectic homology. Preprint. \url{www.mathematik.hu-berlin.de/wendl/pub/SH.pdf}
\bibitem{256}\hypertarget{ref:256}{} E. Witten, Supersymmetry and Morse theory. J. Differential Geometry, 17 (1982), 661–692.
\bibitem{257}\hypertarget{ref:257}{} M. Wong, Dimer models and Hochschild cohomology. arXiv:1908.03005, 2019.
\bibitem{258}\hypertarget{ref:258}{} A. Zimmermann, Representation theory. Algebra and Applications, Springer, Amiens, 2014.
\end{thebibliography}
